% ============================================================
%  PEMBAHASAN SUPER DETAIL — LATIHAN SOAL & STUDI KASUS
%  BAB 4: TRIGONOMETRI DASAR
%  Catatan: Soal Akhir Bab TIDAK dibahas di sini (sesuai permintaan).
% ============================================================
% ============================================================
%  PREAMBUL BERSAMA — BUKU PEMBAHASAN LATIHAN SOAL
%  Dipakai oleh MAE-2026-2027-Latihan1.tex ... Latihan9.tex
%  Kompilasi: pdfLaTeX (disarankan Overleaf / MiKTeX)
% ============================================================
\documentclass[11pt,a4paper]{report}

% ---------- PAKET ----------
\usepackage[utf8]{inputenc}
\usepackage[T1]{fontenc}
\usepackage[bahasa]{babel}
\usepackage{amsmath,amssymb,amsthm,mathtools}
\usepackage{geometry}
\usepackage{xcolor}
\usepackage{graphicx}
\usepackage{enumitem}
\usepackage{tikz}
\usetikzlibrary{calc,arrows.meta,angles,quotes,positioning,decorations.pathreplacing,patterns}
\usepackage{pgfplots}
\pgfplotsset{compat=1.18}
\usepackage[most]{tcolorbox}
\usepackage{fancyhdr}
\usepackage{titlesec}
\usepackage{array,booktabs}
\usepackage{multicol}
\usepackage{pifont}
\usepackage[colorlinks=true,linkcolor=biru,urlcolor=biru]{hyperref}

\geometry{a4paper,margin=2.5cm,headheight=15pt}

% ---------- WARNA ----------
\definecolor{biru}{HTML}{1F4E79}
\definecolor{birumuda}{HTML}{D6E4F0}
\definecolor{hijau}{HTML}{2E7D32}
\definecolor{hijaumuda}{HTML}{E3F2E1}
\definecolor{oranye}{HTML}{C75B12}
\definecolor{oranyemuda}{HTML}{FBE8DA}
\definecolor{abu}{HTML}{F2F2F2}
\definecolor{ungu}{HTML}{6A1B9A}
\definecolor{ungumuda}{HTML}{EDE2F3}
\definecolor{merah}{HTML}{C62828}
\definecolor{merahmuda}{HTML}{FCE4E4}
\definecolor{tosca}{HTML}{00838F}
\definecolor{toscamuda}{HTML}{E0F2F1}
\definecolor{emas}{HTML}{8D6E00}
\definecolor{emasmuda}{HTML}{FFF6D6}

% ---------- HEADER / FOOTER ----------
\newcommand{\babjudul}{Pembahasan Latihan}
\pagestyle{fancy}
\fancyhf{}
\renewcommand{\headrulewidth}{0.8pt}
\renewcommand{\footrulewidth}{0pt}
\fancyhead[L]{\small\textcolor{ungu}{\leftmark}}
\fancyhead[R]{\small\textcolor{ungu}{\babjudul}}
\fancyfoot[C]{\thepage}

% ---------- GAYA JUDUL BAB ----------
\titleformat{\chapter}[display]
  {\normalfont\huge\bfseries\color{ungu}}
  {\filright\large PEMBAHASAN}{8pt}{\Huge}
\titlespacing*{\chapter}{0pt}{0pt}{20pt}

% ---------- LINGKUNGAN TEOREMA ----------
\newtheoremstyle{gaya}{6pt}{6pt}{}{}{\bfseries\color{biru}}{.}{.5em}{}
\theoremstyle{gaya}
\newtheorem{definisi}{Definisi}[chapter]
\newtheorem{sifat}{Sifat}[chapter]

% ---------- KOTAK: SOAL (ungu) ----------
\newtcolorbox{soalbox}[1]{
  enhanced, breakable, colback=ungumuda, colframe=ungu, boxrule=0.8pt,
  arc=3pt, left=10pt, right=10pt, top=8pt, bottom=8pt,
  title={\textbf{Soal #1}}, fonttitle=\bfseries\color{white}, coltitle=white,
  attach boxed title to top left={xshift=10pt,yshift=-10pt},
  boxed title style={colback=ungu,arc=2pt}}

% ---------- KOTAK: KONSEP KUNCI (hijau) ----------
\newtcolorbox{ingat}[1]{
  enhanced, breakable, colback=hijaumuda, colframe=hijau, boxrule=0.6pt,
  arc=3pt, left=10pt, right=10pt, top=6pt, bottom=6pt,
  title={\textbf{Konsep Kunci: #1}}, fonttitle=\bfseries\color{white}, coltitle=white,
  attach boxed title to top left={xshift=10pt,yshift=-10pt},
  boxed title style={colback=hijau,arc=2pt}}

% ---------- KOTAK: JEBAKAN (merah) ----------
\newtcolorbox{jebakan}{
  enhanced, breakable, colback=merahmuda, colframe=merah, boxrule=0.6pt,
  arc=3pt, left=10pt, right=10pt, top=6pt, bottom=6pt,
  title={\textbf{\ding{55}\ Jebakan \& Kesalahan yang Sering Terjadi}},
  fonttitle=\bfseries\color{white}, coltitle=white,
  attach boxed title to top left={xshift=10pt,yshift=-10pt},
  boxed title style={colback=merah,arc=2pt}}

% ---------- KOTAK: VARIASI SOAL (oranye) ----------
\newtcolorbox{variasi}{
  enhanced, breakable, colback=oranyemuda, colframe=oranye, boxrule=0.6pt,
  arc=3pt, left=10pt, right=10pt, top=6pt, bottom=6pt,
  title={\textbf{\ding{72}\ Bentuk Modifikasi Soal Ini}},
  fonttitle=\bfseries\color{white}, coltitle=white,
  attach boxed title to top left={xshift=10pt,yshift=-10pt},
  boxed title style={colback=oranye,arc=2pt}}

% ---------- KOTAK: STUDI KASUS (biru) ----------
\newtcolorbox{studikasus}[1]{
  enhanced, breakable, colback=abu, colframe=biru, boxrule=0.8pt,
  arc=3pt, left=10pt, right=10pt, top=8pt, bottom=8pt,
  title={\textbf{Studi Kasus: #1}}, fonttitle=\bfseries\color{white}, coltitle=white,
  attach boxed title to top left={xshift=10pt,yshift=-10pt},
  boxed title style={colback=biru,arc=2pt}}

% ---------- KOTAK: TIPS (tosca) ----------
\newtcolorbox{tips}{
  enhanced, breakable, colback=toscamuda, colframe=tosca, boxrule=0.6pt,
  arc=3pt, left=10pt, right=10pt, top=6pt, bottom=6pt,
  borderline west={3pt}{0pt}{tosca}}

% ---------- PERINTAH BANTU ----------
\newcommand{\jawab}{\par\smallskip\noindent\textit{\textcolor{oranye}{Penyelesaian langkah demi langkah:}}\par\smallskip}
\newcommand{\langkah}[1]{\par\smallskip\noindent\textbf{\textcolor{biru}{Langkah #1.}}\ }
\newcommand{\jadi}{\par\smallskip\noindent$\Rightarrow$\ \textbf{Jadi, }}
% Judul tiap nomor soal
\newcommand{\soalno}[2]{%
  \par\bigskip
  \noindent{\large\bfseries\color{ungu}\rule[-2pt]{4pt}{14pt}\ Soal #1}%
  \ {\normalfont\itshape\color{gray}(#2)}\par\nopagebreak\smallskip}

% ---------- HALAMAN JUDUL OTOMATIS ----------
% #1 = judul topik (mis. "Barisan dan Deret"); #2 = nomor bab (angka)
\newcommand{\buathalamanjudul}[2]{%
  \begin{titlepage}
  \centering
  \vspace*{2cm}
  {\color{ungu}\rule{\textwidth}{2pt}}\\[0.6cm]
  {\Huge\bfseries\color{ungu} PEMBAHASAN\\[0.3cm] LATIHAN SOAL}\\[0.4cm]
  {\Large BAB #2 --- #1}\\[0.2cm]
  {\color{ungu}\rule{\textwidth}{2pt}}\\[1.2cm]

  {\large Matematika SMA --- Fase E (Kelas 10)}\\[0.2cm]
  {\large Kurikulum Merdeka}\\[2cm]

  \begin{tcolorbox}[enhanced,width=0.85\textwidth,colback=ungumuda,colframe=ungu,
    arc=4pt,boxrule=1pt]
  \centering\small
  Buku pembahasan ini mengupas \textbf{setiap soal pada kotak ungu (Latihan Soal)}
  di tiap subbab, \emph{selangkah demi selangkah seperti mengajari dari nol}:
  konsep yang dipakai, jebakan yang sering menjerat, cara soal bisa dimodifikasi,
  dan visualisasi. \emph{Soal Akhir Bab tidak dibahas di sini.}
  \end{tcolorbox}

  \vfill
  {\large Tahun Pelajaran 2026/2027}
  \end{titlepage}
  \tableofcontents
  \thispagestyle{empty}
  \clearpage}

% ============================================================
%  PENJAGA KOMPILASI
%  File ini adalah PREAMBUL BERSAMA (di-\input oleh Latihan1..9),
%  BUKAN untuk dikompilasi sendiri. Jika editor terlanjur mengompilasi
%  file ini secara langsung, tampilkan satu halaman info (bukan error).
%  Saat di-\input oleh berkas Latihan, blok ini dilewati.
% ============================================================
\ifnum\pdfstrcmp{\jobname}{MAE-2026-2027-Latihan-preamble}=0
  \begin{document}
  \thispagestyle{empty}
  \begin{center}
  \vspace*{3cm}
  {\Huge\bfseries\color{ungu} Berkas Preambel Bersama}\\[1cm]
  \begin{tcolorbox}[enhanced,width=0.85\textwidth,colback=ungumuda,
    colframe=ungu,arc=4pt,boxrule=1pt]
  \centering
  File ini \textbf{bukan untuk dikompilasi sendiri.} Ia berisi pengaturan
  bersama (paket, warna, kotak, perintah) yang dipanggil lewat
  \texttt{\textbackslash input} oleh berkas \textbf{MAE-2026-2027-Latihan1.tex}
  sampai \textbf{Latihan9.tex}.

  \medskip
  Silakan buka dan kompilasi salah satu berkas \textbf{Latihan1--Latihan9}
  tersebut untuk menghasilkan PDF pembahasan.
  \end{tcolorbox}
  \end{center}
  \end{document}
\fi

\renewcommand{\babjudul}{Pembahasan Latihan — Bab 4}

\begin{document}
\setcounter{chapter}{4}
\buathalamanjudul{Trigonometri Dasar}{4}

% ============================================================
\chapter*{Pembahasan Bab 4: Trigonometri Dasar}
\addcontentsline{toc}{chapter}{Pembahasan Bab 4: Trigonometri Dasar}
\markboth{PEMBAHASAN BAB 4}{}
% ============================================================

\begin{tips}
\textbf{Cara membaca buku ini.} Tiap soal: \textbf{(1)} soal ditulis ulang,
\textbf{(2)} konsep kunci, \textbf{(3)} langkah demi langkah, \textbf{(4)}
jebakan, \textbf{(5)} variasi. \emph{Aturan emas Bab 4:} gambar segitiganya,
tandai sudut acuan lebih dulu, baru tentukan mana ``depan/samping/miring''.
\end{tips}

% ============================================================
\section*{Studi Kasus: Mengukur Tinggi Menara}
\addcontentsline{toc}{section}{Studi Kasus: Mengukur Tinggi Menara}
% ============================================================

\begin{studikasus}{Mengukur Tinggi Menara}
Seorang siswa berdiri $50$ m dari kaki menara. Dengan klinometer ia mengukur
sudut elevasi puncak menara $40^\circ$. Tinggi matanya $1{,}5$ m.
\begin{center}
\begin{tikzpicture}[scale=0.85]
  \coordinate (A) at (0,0); \coordinate (B) at (5,0); \coordinate (C) at (5,4);
  \draw[very thick,biru] (A)--(B)--(C)--cycle;
  \draw (4.7,0) rectangle (5,0.3);
  \pic[draw,angle radius=0.9cm,"$40^\circ$"] {angle=B--A--C} ;
  \node[below] at (2.5,0) {50 m};
  \node[right] at (5,2) {tinggi $=?$};
  \fill (0,0) circle (2pt) node[left]{mata};
\end{tikzpicture}
\end{center}
\textbf{Pertanyaan.} (a) Besaran apa yang diketahui dan apa yang dicari?
(b) Perbandingan trigonometri mana yang menghubungkan keduanya?
(c) Bagaimana menambahkan tinggi mata? ($\tan40^\circ\approx0{,}839$)
\end{studikasus}

\subsection*{Jawaban lengkap studi kasus}

\begin{ingat}{Memilih sin/cos/tan}
Pilih perbandingan dari \emph{sisi yang diketahui} dan \emph{sisi yang dicari}
relatif terhadap sudut: depan/miring $\to\sin$; samping/miring $\to\cos$;
depan/samping $\to\tan$.
\end{ingat}

\textbf{(a) Yang diketahui vs dicari.} Diketahui jarak mendatar ke menara $=50$ m
(ini sisi \emph{samping} sudut $40^\circ$). Dicari tinggi menara di atas mata
(ini sisi \emph{depan} sudut). Sisi miring tidak terlibat.

\textbf{(b) Perbandingan yang tepat.} Karena melibatkan \emph{depan} dan
\emph{samping}, gunakan $\tan$:
\[
\tan40^\circ=\frac{\text{depan}}{\text{samping}}=\frac{h}{50}
\;\Rightarrow\; h=50\tan40^\circ\approx50(0{,}839)=41{,}95\ \text{m}.
\]

\textbf{(c) Menambahkan tinggi mata.} Nilai $h=41{,}95$ m dihitung dari
\emph{ketinggian mata}, bukan dari tanah. Maka tinggi menara sebenarnya:
\[
41{,}95 + 1{,}5 \approx \textbf{43{,}5 m}.
\]
Inilah kekuatan trigonometri: mengubah sebuah \emph{sudut} menjadi jarak yang
tak terjangkau tanpa memanjat.

% ############################################################
\section*{Subbab 4.1 — Perbandingan Trigonometri}
\addcontentsline{toc}{section}{Subbab 4.1 — Perbandingan Trigonometri}
% ############################################################

\begin{ingat}{SOH-CAH-TOA \& sudut istimewa}
\[
\sin\alpha=\frac{\text{de}}{\text{mi}},\quad
\cos\alpha=\frac{\text{sa}}{\text{mi}},\quad
\tan\alpha=\frac{\text{de}}{\text{sa}}.
\]
\[
\begin{array}{c|ccccc}
\alpha & 0^\circ & 30^\circ & 45^\circ & 60^\circ & 90^\circ\\\hline
\sin & 0 & \frac12 & \frac{\sqrt2}{2} & \frac{\sqrt3}{2} & 1\\[3pt]
\cos & 1 & \frac{\sqrt3}{2} & \frac{\sqrt2}{2} & \frac12 & 0\\[3pt]
\tan & 0 & \frac{\sqrt3}{3} & 1 & \sqrt3 & -\\
\end{array}
\]
Identitas pokok: $\sin^2\alpha+\cos^2\alpha=1$.
\end{ingat}

% ---------------- SOAL 4.1.1 ----------------
\soalno{4.1.1}{tiga rasio dari dua sisi}
\begin{soalbox}{}
Diketahui sisi depan $=6$ dan sisi samping $=8$. Tentukan $\sin\alpha$,
$\cos\alpha$, dan $\tan\alpha$.
\end{soalbox}

\jawab
\langkah{1} Cari sisi miring dengan Pythagoras:
$\text{mi}=\sqrt{6^2+8^2}=\sqrt{36+64}=\sqrt{100}=10$.
\langkah{2} Terapkan definisi:
\[
\sin\alpha=\frac{6}{10}=\frac35,\quad
\cos\alpha=\frac{8}{10}=\frac45,\quad
\tan\alpha=\frac{6}{8}=\frac34.
\]
\jadi $\sin\alpha=\frac35,\ \cos\alpha=\frac45,\ \tan\alpha=\frac34$
(segitiga $3$-$4$-$5$ yang terkenal).

\begin{jebakan}
\begin{itemize}[leftmargin=*,topsep=2pt]
  \item \textbf{Lupa mencari miring dulu.} $\sin$ dan $\cos$ butuh sisi miring
  ($=10$), bukan langsung memakai $6$ dan $8$ sebagai penyebut.
  \item \textbf{Tertukar $\sin$ dan $\cos$.} $\sin$ pakai sisi \emph{depan};
  $\cos$ pakai sisi \emph{samping}.
\end{itemize}
\end{jebakan}

\begin{variasi}
Diberi depan \& miring (cari samping dulu), atau segitiga $5$-$12$-$13$. Cek
selalu $\sin^2+\cos^2=1$: $\frac{9}{25}+\frac{16}{25}=1$ \checkmark.
\end{variasi}

% ---------------- SOAL 4.1.2 ----------------
\soalno{4.1.2}{tangga bersandar (sisi depan)}
\begin{soalbox}{}
Sebuah tangga bersandar membentuk sudut $60^\circ$ dengan tanah. Jika panjang
tangga $4$ m, berapa tinggi dinding yang dicapai?
\end{soalbox}

\jawab
\langkah{1} Identifikasi: panjang tangga = sisi \emph{miring} ($4$ m); tinggi
dinding = sisi \emph{depan} sudut $60^\circ$.
\langkah{2} Depan \& miring $\Rightarrow$ gunakan $\sin$:
\[
\sin60^\circ=\frac{\text{tinggi}}{4}\Rightarrow \text{tinggi}=4\sin60^\circ
=4\cdot\frac{\sqrt3}{2}=2\sqrt3\approx3{,}46\ \text{m}.
\]
\jadi tinggi dinding $=2\sqrt3\approx3{,}46$ m.

\begin{jebakan}
\begin{itemize}[leftmargin=*,topsep=2pt]
  \item \textbf{Memakai $\cos$.} $\cos60^\circ$ memberi jarak \emph{kaki tangga
  ke dinding} (sisi samping), bukan tinggi. Baca yang ditanya.
  \item \textbf{$\sin60^\circ=\frac12$?} Bukan; $\sin60^\circ=\frac{\sqrt3}{2}$.
\end{itemize}
\end{jebakan}

\begin{variasi}
Ditanya jarak kaki tangga ke dinding $\Rightarrow \cos60^\circ\cdot4=2$ m.
Sudut lain ($30^\circ$, $45^\circ$) mengubah nilai tapi metode sama.
\end{variasi}

% ---------------- SOAL 4.1.3 ----------------
\soalno{4.1.3}{penjumlahan nilai istimewa}
\begin{soalbox}{}
Hitung $\sin45^\circ+\cos45^\circ$.
\end{soalbox}

\jawab
\langkah{1} $\sin45^\circ=\dfrac{\sqrt2}{2}$ dan $\cos45^\circ=\dfrac{\sqrt2}{2}$
(pada segitiga siku-siku sama kaki keduanya sama).
\langkah{2} Jumlahkan: $\dfrac{\sqrt2}{2}+\dfrac{\sqrt2}{2}=\dfrac{2\sqrt2}{2}
=\sqrt2$.
\jadi $\sin45^\circ+\cos45^\circ=\sqrt2\approx1{,}41$.

\begin{jebakan}
\begin{itemize}[leftmargin=*,topsep=2pt]
  \item \textbf{Mengira hasilnya $1$.} $\sin45+\cos45\neq\sin^2+\cos^2$.
  Jumlah nilainya $\sqrt2$, bukan $1$.
\end{itemize}
\end{jebakan}

\begin{variasi}
$\sin30^\circ+\cos60^\circ=1$; $\sin60^\circ-\cos30^\circ=0$ (keduanya
$\frac{\sqrt3}{2}$).
\end{variasi}

% ---------------- SOAL 4.1.4 ----------------
\soalno{4.1.4}{identitas Pythagoras}
\begin{soalbox}{}
Jika $\cos\alpha=\tfrac{12}{13}$ ($\alpha$ lancip), tentukan $\sin\alpha$.
\end{soalbox}

\begin{ingat}{Identitas $\sin^2+\cos^2=1$}
Ini Pythagoras yang menyamar. Jika satu rasio diketahui, yang lain dapat dicari
tanpa mengetahui sudutnya. Untuk $\alpha$ lancip (kuadran I), semua rasio positif.
\end{ingat}

\jawab
\langkah{1} Dari $\sin^2\alpha=1-\cos^2\alpha$:
\[
\sin^2\alpha=1-\left(\frac{12}{13}\right)^2=1-\frac{144}{169}=\frac{25}{169}.
\]
\langkah{2} Akarkan (ambil positif karena lancip):
$\sin\alpha=\dfrac{5}{13}$.
\jadi $\sin\alpha=\dfrac{5}{13}$ (segitiga $5$-$12$-$13$).

\begin{jebakan}
\begin{itemize}[leftmargin=*,topsep=2pt]
  \item \textbf{Lupa mengkuadratkan.} $\sin\alpha\neq1-\frac{12}{13}$. Yang
  dikurangi adalah \emph{kuadratnya}.
  \item \textbf{Salah tanda di kuadran lain.} Kalau $\alpha$ tumpul (kuadran II),
  $\sin$ tetap $+$ tapi cek konteks. Di sini lancip $\Rightarrow +$.
\end{itemize}
\end{jebakan}

\begin{variasi}
Diberi $\sin$ cari $\cos$; atau diberi $\tan$ cari $\sin,\cos$ (bentuk segitiga
dari $\tan=\frac{\text{de}}{\text{sa}}$).
\end{variasi}

% ---------------- SOAL 4.1.5 ----------------
\soalno{4.1.5}{hasil kali tan istimewa}
\begin{soalbox}{}
Tentukan nilai $\tan60^\circ\cdot\tan30^\circ$.
\end{soalbox}

\jawab
\langkah{1} $\tan60^\circ=\sqrt3$ dan $\tan30^\circ=\dfrac{1}{\sqrt3}$
$\left(=\dfrac{\sqrt3}{3}\right)$.
\langkah{2} Kalikan: $\sqrt3\cdot\dfrac{1}{\sqrt3}=1$.
\jadi $\tan60^\circ\cdot\tan30^\circ=1$.

\begin{tips}
Ini bukan kebetulan: $30^\circ$ dan $60^\circ$ saling berpenyiku
($30+60=90$), dan $\tan\alpha\cdot\tan(90^\circ-\alpha)=1$ selalu, sebab
$\tan(90^\circ-\alpha)=\cot\alpha=\frac{1}{\tan\alpha}$.
\end{tips}

\begin{jebakan}
\begin{itemize}[leftmargin=*,topsep=2pt]
  \item \textbf{$\tan30^\circ=\sqrt3$?} Terbalik. $\tan30^\circ=\frac{1}{\sqrt3}$,
  $\tan60^\circ=\sqrt3$.
\end{itemize}
\end{jebakan}

\begin{variasi}
$\sin\alpha\cdot\csc\alpha=1$; $\tan45^\circ\cdot\tan45^\circ=1$;
$\tan20^\circ\cdot\tan70^\circ=1$.
\end{variasi}

% ---------------- SOAL 4.1.6 ----------------
\soalno{4.1.6}{layang-layang (sisi depan)}
\begin{soalbox}{}
Sebuah layang-layang terbang dengan benang $100$ m membentuk sudut $30^\circ$
terhadap tanah. Berapa tingginya?
\end{soalbox}

\jawab
\langkah{1} Benang = sisi miring ($100$ m); tinggi = sisi depan sudut $30^\circ$.
\langkah{2} $\sin30^\circ=\dfrac{\text{tinggi}}{100}\Rightarrow
\text{tinggi}=100\sin30^\circ=100\cdot\tfrac12=50$ m.
\jadi tinggi layang-layang $=50$ m.

\begin{jebakan}
\begin{itemize}[leftmargin=*,topsep=2pt]
  \item \textbf{Menganggap benang lurus vertikal.} Benang miring; tinggi hanya
  komponen vertikalnya ($\sin$).
\end{itemize}
\end{jebakan}

\begin{variasi}
Ditanya jarak mendatar layang-layang $\Rightarrow \cos30^\circ\cdot100
=50\sqrt3$ m.
\end{variasi}

% ---------------- SOAL 4.1.7 ----------------
\soalno{4.1.7}{membuktikan $\tan=\sin/\cos$}
\begin{soalbox}{}
Buktikan $\dfrac{\sin\alpha}{\cos\alpha}=\tan\alpha$ dari definisi.
\end{soalbox}

\jawab
\langkah{1} Tulis $\sin$ dan $\cos$ dengan definisi rasio:
\[
\frac{\sin\alpha}{\cos\alpha}=\frac{\ \text{de}/\text{mi}\ }{\ \text{sa}/\text{mi}\ }.
\]
\langkah{2} Bagi pecahan = kali kebalikan; ``mi'' saling coret:
\[
=\frac{\text{de}}{\text{mi}}\cdot\frac{\text{mi}}{\text{sa}}
=\frac{\text{de}}{\text{sa}}=\tan\alpha.
\]
\jadi terbukti $\dfrac{\sin\alpha}{\cos\alpha}=\tan\alpha$. $\blacksquare$

\begin{jebakan}
\begin{itemize}[leftmargin=*,topsep=2pt]
  \item \textbf{Membuktikan dengan satu contoh angka.} Bukti sejati memakai
  definisi umum, bukan sekadar mengecek $\alpha=30^\circ$.
\end{itemize}
\end{jebakan}

\begin{variasi}
Buktikan $\cot\alpha=\frac{\cos\alpha}{\sin\alpha}$; atau
$\sec\alpha=\frac{1}{\cos\alpha}$.
\end{variasi}

% ---------------- SOAL 4.1.8 ----------------
\soalno{4.1.8}{identitas pokok bernilai 1}
\begin{soalbox}{}
Hitung $\sin^2 30^\circ+\cos^2 30^\circ$.
\end{soalbox}

\jawab
\textbf{Cara cepat (identitas).} $\sin^2\alpha+\cos^2\alpha=1$ untuk sudut
\emph{apa pun}, jadi langsung $=1$.

\textbf{Cara panjang (bukti).}
$\left(\tfrac12\right)^2+\left(\tfrac{\sqrt3}{2}\right)^2
=\tfrac14+\tfrac34=1$. \checkmark
\jadi hasilnya $1$.

\begin{jebakan}
\begin{itemize}[leftmargin=*,topsep=2pt]
  \item \textbf{$\sin^2\alpha$ ditulis $\sin\alpha^2$.} $\sin^2\alpha=(\sin\alpha)^2$,
  bukan $\sin(\alpha^2)$.
\end{itemize}
\end{jebakan}

\begin{variasi}
Berlaku untuk sudut mana pun: $\sin^2 17^\circ+\cos^2 17^\circ=1$. Digunakan
untuk menyederhanakan ekspresi rumit (lihat Soal 4.2.9).
\end{variasi}

% ---------------- SOAL 4.1.9 ----------------
\soalno{4.1.9}{bayangan \& tinggi (tan)}
\begin{soalbox}{}
Sudut elevasi matahari $60^\circ$, panjang bayangan sebuah tiang $3$ m. Berapa
tinggi tiang?
\end{soalbox}

\jawab
\langkah{1} Bayangan = sisi samping (mendatar); tinggi tiang = sisi depan sudut
$60^\circ$. Depan \& samping $\Rightarrow\tan$.
\langkah{2} $\tan60^\circ=\dfrac{\text{tinggi}}{3}\Rightarrow
\text{tinggi}=3\tan60^\circ=3\sqrt3\approx5{,}20$ m.
\jadi tinggi tiang $=3\sqrt3\approx5{,}20$ m.

\begin{jebakan}
\begin{itemize}[leftmargin=*,topsep=2pt]
  \item \textbf{Membalik: $\text{tinggi}=3/\tan60^\circ$.} Karena tinggi adalah
  ``depan'' dan bayangan ``samping'', tinggi $=$ samping $\times\tan$, jadi
  \emph{dikali}, bukan dibagi.
\end{itemize}
\end{jebakan}

\begin{variasi}
Sudut elevasi kecil $\Rightarrow$ bayangan panjang. ``Berapa panjang bayangan
jika tinggi tiang $6$ m dan elevasi $30^\circ$?''
\end{variasi}

% ---------------- SOAL 4.1.10 ----------------
\soalno{4.1.10}{dari nilai tan ke sudut}
\begin{soalbox}{}
Jika $\tan\alpha=1$, tentukan $\alpha$ (lancip) dan $\sin\alpha$.
\end{soalbox}

\jawab
\langkah{1} $\tan\alpha=1$ berarti sisi depan $=$ samping (segitiga siku-siku
sama kaki). Sudut lancip yang memenuhinya $\alpha=45^\circ$.
\langkah{2} $\sin45^\circ=\dfrac{\sqrt2}{2}$.
\jadi $\alpha=45^\circ$ dan $\sin\alpha=\dfrac{\sqrt2}{2}$.

\begin{jebakan}
\begin{itemize}[leftmargin=*,topsep=2pt]
  \item \textbf{Lupa batasan ``lancip''.} $\tan\alpha=1$ juga di $225^\circ$;
  batasan lancip menyisakan $45^\circ$ saja.
\end{itemize}
\end{jebakan}

\begin{variasi}
$\tan\alpha=\sqrt3\Rightarrow\alpha=60^\circ$; $\tan\alpha=\frac{1}{\sqrt3}
\Rightarrow\alpha=30^\circ$.
\end{variasi}

% ############################################################
\section*{Subbab 4.2 — Sudut di Berbagai Kuadran dan Identitas}
\addcontentsline{toc}{section}{Subbab 4.2 — Kuadran \& Identitas}
% ############################################################

\begin{ingat}{ASTC \& sudut acuan}
Di kuadran I \textbf{A}ll positif; II hanya \textbf{S}in; III hanya \textbf{T}an;
IV hanya \textbf{C}os. Nilai mutlaknya diambil dari \emph{sudut acuan} (jarak ke
sumbu $x$ terdekat), lalu tanda ditentukan ASTC. Sudut berelasi:
\[
\sin(90^\circ-\alpha)=\cos\alpha,\quad 1+\tan^2\alpha=\sec^2\alpha.
\]
\end{ingat}

% ---------------- SOAL 4.2.1 ----------------
\soalno{4.2.1}{nilai di kuadran II}
\begin{soalbox}{}
Tentukan $\sin120^\circ$, $\cos120^\circ$, dan $\tan120^\circ$.
\end{soalbox}

\jawab
\langkah{1} Tentukan kuadran: $120^\circ$ di kuadran II ($90^\circ$--$180^\circ$).
Di sini hanya \textbf{Sin} positif.
\langkah{2} Sudut acuan: $180^\circ-120^\circ=60^\circ$.
\langkah{3} Ambil nilai $60^\circ$ lalu beri tanda ASTC:
\[
\sin120^\circ=+\sin60^\circ=\frac{\sqrt3}{2},\quad
\cos120^\circ=-\cos60^\circ=-\frac12,\quad
\tan120^\circ=-\tan60^\circ=-\sqrt3.
\]
\jadi $\sin120^\circ=\frac{\sqrt3}{2},\ \cos120^\circ=-\frac12,\ \tan120^\circ=-\sqrt3$.

\begin{jebakan}
\begin{itemize}[leftmargin=*,topsep=2pt]
  \item \textbf{Sudut acuan salah.} Di kuadran II acuan $=180^\circ-\alpha$
  (bukan $\alpha-90^\circ$).
  \item \textbf{Lupa tanda.} $\cos$ dan $\tan$ negatif di kuadran II.
\end{itemize}
\end{jebakan}

\begin{variasi}
$\sin210^\circ$ (KIII, acuan $30^\circ$, sin$-$); $\cos330^\circ$ (KIV, acuan
$30^\circ$, cos$+$).
\end{variasi}

% ---------------- SOAL 4.2.2 ----------------
\soalno{4.2.2}{nilai di sumbu}
\begin{soalbox}{}
Hitung $\cos180^\circ+\sin90^\circ$.
\end{soalbox}

\jawab
\langkah{1} Nilai di sumbu (lingkaran satuan): $\cos180^\circ=-1$ (titik
$(-1,0)$); $\sin90^\circ=1$ (titik $(0,1)$).
\langkah{2} Jumlahkan: $-1+1=0$.
\jadi $\cos180^\circ+\sin90^\circ=0$.

\begin{jebakan}
\begin{itemize}[leftmargin=*,topsep=2pt]
  \item \textbf{$\cos180^\circ=1$?} Bukan; titik di $180^\circ$ adalah $(-1,0)$,
  jadi $\cos=-1$. Bayangkan lingkaran satuan.
\end{itemize}
\end{jebakan}

\begin{variasi}
$\sin180^\circ=0$, $\cos90^\circ=0$, $\sin270^\circ=-1$. Hafalkan lewat
koordinat lingkaran satuan, bukan hafalan buta.
\end{variasi}

% ---------------- SOAL 4.2.3 ----------------
\soalno{4.2.3}{identitas dengan tanda kuadran}
\begin{soalbox}{}
Jika $\sin\alpha=\tfrac45$ dan $\alpha$ di kuadran II, tentukan $\cos\alpha$.
\end{soalbox}

\jawab
\langkah{1} $\cos^2\alpha=1-\sin^2\alpha=1-\frac{16}{25}=\frac{9}{25}$, jadi
$\cos\alpha=\pm\frac35$.
\langkah{2} \textbf{Pilih tanda dari kuadran.} Di kuadran II, $\cos$
\emph{negatif}. Maka $\cos\alpha=-\frac35$.
\jadi $\cos\alpha=-\dfrac35$.

\begin{jebakan}
\begin{itemize}[leftmargin=*,topsep=2pt]
  \item \textbf{Mengambil $+\frac35$ otomatis.} Identitas memberi $\pm$; kuadran
  yang menentukan tanda. Ini kesalahan paling sering.
\end{itemize}
\end{jebakan}

\begin{variasi}
$\cos\alpha$ diberi, $\alpha$ di KIII $\Rightarrow\sin$ negatif. Selalu tanya:
di kuadran ini, tanda fungsinya apa?
\end{variasi}

% ---------------- SOAL 4.2.4 ----------------
\soalno{4.2.4}{sudut berelasi + kebalikan}
\begin{soalbox}{}
Sederhanakan $\sin(90^\circ-\alpha)\cdot\sec\alpha$.
\end{soalbox}

\jawab
\langkah{1} $\sin(90^\circ-\alpha)=\cos\alpha$ (sudut berpenyiku).
\langkah{2} $\sec\alpha=\dfrac{1}{\cos\alpha}$.
\langkah{3} Kalikan: $\cos\alpha\cdot\dfrac{1}{\cos\alpha}=1$.
\jadi hasilnya $1$.

\begin{jebakan}
\begin{itemize}[leftmargin=*,topsep=2pt]
  \item \textbf{$\sin(90^\circ-\alpha)=\sin90^\circ-\sin\alpha$?} Salah! Fungsi
  trigonometri tidak menembus tanda kurang. Yang benar $=\cos\alpha$.
\end{itemize}
\end{jebakan}

\begin{variasi}
$\cos(90^\circ-\alpha)\csc\alpha=1$; $\tan(90^\circ-\alpha)=\cot\alpha$.
\end{variasi}

% ---------------- SOAL 4.2.5 ----------------
\soalno{4.2.5}{menurunkan identitas}
\begin{soalbox}{}
Buktikan $1+\tan^2\alpha=\sec^2\alpha$ dari identitas pokok.
\end{soalbox}

\jawab
\langkah{1} Mulai dari identitas pokok $\sin^2\alpha+\cos^2\alpha=1$.
\langkah{2} Bagi \emph{seluruhnya} dengan $\cos^2\alpha$:
\[
\frac{\sin^2\alpha}{\cos^2\alpha}+\frac{\cos^2\alpha}{\cos^2\alpha}
=\frac{1}{\cos^2\alpha}.
\]
\langkah{3} Kenali tiap suku: $\frac{\sin^2}{\cos^2}=\tan^2\alpha$;
$\frac{\cos^2}{\cos^2}=1$; $\frac{1}{\cos^2}=\sec^2\alpha$. Maka:
\[
\tan^2\alpha+1=\sec^2\alpha. \quad\blacksquare
\]

\begin{jebakan}
\begin{itemize}[leftmargin=*,topsep=2pt]
  \item \textbf{Membagi sebagian saja.} Semua tiga suku harus dibagi
  $\cos^2\alpha$, termasuk ruas kanan.
\end{itemize}
\end{jebakan}

\begin{variasi}
Bagi dengan $\sin^2\alpha$ untuk memperoleh $1+\cot^2\alpha=\csc^2\alpha$.
\end{variasi}

% ---------------- SOAL 4.2.6 ----------------
\soalno{4.2.6}{nilai di kuadran III}
\begin{soalbox}{}
Tentukan $\tan225^\circ$.
\end{soalbox}

\jawab
\langkah{1} $225^\circ$ di kuadran III ($180^\circ$--$270^\circ$); di sini
\textbf{Tan} positif.
\langkah{2} Sudut acuan: $225^\circ-180^\circ=45^\circ$.
\langkah{3} $\tan225^\circ=+\tan45^\circ=1$.
\jadi $\tan225^\circ=1$.

\begin{jebakan}
\begin{itemize}[leftmargin=*,topsep=2pt]
  \item \textbf{Memberi tanda negatif.} $\tan$ justru \emph{positif} di kuadran
  III (sin dan cos sama-sama negatif, bagi keduanya jadi positif).
  \item \textbf{Sudut acuan di KIII} $=\alpha-180^\circ$.
\end{itemize}
\end{jebakan}

\begin{variasi}
$\sin225^\circ=-\frac{\sqrt2}{2}$, $\cos225^\circ=-\frac{\sqrt2}{2}$;
cek $\tan=\frac{\sin}{\cos}=1$ \checkmark.
\end{variasi}

% ---------------- SOAL 4.2.7 ----------------
\soalno{4.2.7}{jumlah kuadrat sin}
\begin{soalbox}{}
Hitung $\sin^2 60^\circ+\sin^2 30^\circ$.
\end{soalbox}

\jawab
\langkah{1} $\sin60^\circ=\frac{\sqrt3}{2}\Rightarrow\sin^2 60^\circ=\frac34$;
$\sin30^\circ=\frac12\Rightarrow\sin^2 30^\circ=\frac14$.
\langkah{2} Jumlahkan: $\frac34+\frac14=1$.
\jadi hasilnya $1$.

\begin{tips}
Bukan kebetulan: karena $\sin30^\circ=\cos60^\circ$, penjumlahan ini setara
$\sin^2 60^\circ+\cos^2 60^\circ=1$ (identitas pokok).
\end{tips}

\begin{jebakan}
\begin{itemize}[leftmargin=*,topsep=2pt]
  \item \textbf{Tertukar dengan $\sin(60+30)$.} Ini jumlah \emph{kuadrat nilai},
  bukan sinus dari jumlah sudut.
\end{itemize}
\end{jebakan}

\begin{variasi}
$\cos^2 45^\circ+\sin^2 45^\circ=1$; $\sin^2 60^\circ-\cos^2 60^\circ
=\frac34-\frac14=\frac12$ (ini $\cos2\cdot? $ ... pengayaan).
\end{variasi}

% ---------------- SOAL 4.2.8 ----------------
\soalno{4.2.8}{menentukan kuadran}
\begin{soalbox}{}
Tentukan kuadran $\alpha$ jika $\sin\alpha>0$ dan $\cos\alpha<0$.
\end{soalbox}

\jawab
\langkah{1} $\sin\alpha>0$: kuadran I atau II (sin positif di atas sumbu $x$).
\langkah{2} $\cos\alpha<0$: kuadran II atau III (cos negatif di kiri sumbu $y$).
\langkah{3} Irisan kedua syarat: kuadran \textbf{II}.
\jadi $\alpha$ di kuadran II.

\begin{jebakan}
\begin{itemize}[leftmargin=*,topsep=2pt]
  \item \textbf{Menebak tanpa mengiris.} Gabungkan \emph{kedua} syarat; hanya
  satu kuadran memenuhi keduanya.
\end{itemize}
\end{jebakan}

\begin{variasi}
$\sin<0$ \& $\tan>0\Rightarrow$ KIII; $\cos>0$ \& $\tan<0\Rightarrow$ KIV.
\end{variasi}

% ---------------- SOAL 4.2.9 ----------------
\soalno{4.2.9}{menyederhanakan lewat identitas}
\begin{soalbox}{}
Sederhanakan $\dfrac{1-\cos^2\alpha}{\sin\alpha}$.
\end{soalbox}

\jawab
\langkah{1} Kenali $1-\cos^2\alpha=\sin^2\alpha$ (dari identitas pokok, disusun
ulang).
\langkah{2} Substitusi:
\[
\frac{1-\cos^2\alpha}{\sin\alpha}=\frac{\sin^2\alpha}{\sin\alpha}=\sin\alpha.
\]
\jadi hasilnya $\sin\alpha$.

\begin{jebakan}
\begin{itemize}[leftmargin=*,topsep=2pt]
  \item \textbf{Mencoba mencoret $\cos^2$ dengan $\sin$.} Tidak bisa langsung;
  ubah dulu pembilang jadi $\sin^2\alpha$ memakai identitas.
\end{itemize}
\end{jebakan}

\begin{variasi}
$\frac{1-\sin^2\alpha}{\cos\alpha}=\cos\alpha$;
$\frac{\sin^2\alpha}{1-\cos\alpha}=1+\cos\alpha$ (pakai selisih kuadrat).
\end{variasi}

% ---------------- SOAL 4.2.10 ----------------
\soalno{4.2.10}{nilai di kuadran IV}
\begin{soalbox}{}
Hitung $\cos300^\circ$.
\end{soalbox}

\jawab
\langkah{1} $300^\circ$ di kuadran IV ($270^\circ$--$360^\circ$); di sini
\textbf{Cos} positif.
\langkah{2} Sudut acuan: $360^\circ-300^\circ=60^\circ$.
\langkah{3} $\cos300^\circ=+\cos60^\circ=\dfrac12$.
\jadi $\cos300^\circ=\dfrac12$.

\begin{jebakan}
\begin{itemize}[leftmargin=*,topsep=2pt]
  \item \textbf{Acuan di KIV} $=360^\circ-\alpha$. Jangan pakai $\alpha-270^\circ$.
\end{itemize}
\end{jebakan}

\begin{variasi}
$\sin300^\circ=-\frac{\sqrt3}{2}$ (sin$-$ di KIV); $\tan300^\circ=-\sqrt3$.
\end{variasi}

% ############################################################
\section*{Subbab 4.3 — Trigonometri Lanjutan Dasar (Pengayaan)}
\addcontentsline{toc}{section}{Subbab 4.3 — Trigonometri Lanjutan (Pengayaan)}
% ############################################################

\begin{ingat}{Rumus majemuk, rangkap, setengah, persamaan}
\[
\sin(A\pm B)=\sin A\cos B\pm\cos A\sin B,\quad
\cos(A\pm B)=\cos A\cos B\mp\sin A\sin B,
\]
\[
\tan(A\pm B)=\frac{\tan A\pm\tan B}{1\mp\tan A\tan B},\quad
\sin2A=2\sin A\cos A,\quad \cos2A=1-2\sin^2 A.
\]
Persamaan: $\sin x=\sin\alpha\Rightarrow x=\alpha$ atau $180^\circ-\alpha$;\quad
$\cos x=\cos\alpha\Rightarrow x=\alpha$ atau $360^\circ-\alpha$;\quad
$\tan x=\tan\alpha\Rightarrow x=\alpha$ atau $\alpha+180^\circ$.
\end{ingat}

% ---------------- SOAL 4.3.1 ----------------
\soalno{4.3.1}{cos sudut campuran}
\begin{soalbox}{}
Hitung $\cos75^\circ$ dengan rumus jumlah sudut.
\end{soalbox}

\jawab
\langkah{1} Pecah $75^\circ=45^\circ+30^\circ$.
\langkah{2} $\cos(A+B)=\cos A\cos B-\sin A\sin B$ (perhatikan tanda
\emph{berlawanan} pada cos):
\[
\cos75^\circ=\cos45^\circ\cos30^\circ-\sin45^\circ\sin30^\circ
=\frac{\sqrt2}{2}\cdot\frac{\sqrt3}{2}-\frac{\sqrt2}{2}\cdot\frac12.
\]
\langkah{3} $=\dfrac{\sqrt6}{4}-\dfrac{\sqrt2}{4}=\dfrac{\sqrt6-\sqrt2}{4}$.
\jadi $\cos75^\circ=\dfrac{\sqrt6-\sqrt2}{4}\approx0{,}259$.

\begin{jebakan}
\begin{itemize}[leftmargin=*,topsep=2pt]
  \item \textbf{Tanda pada cos.} Rumus $\cos(A+B)$ memakai tanda \emph{minus}
  di tengah (kebalikan dari sin). Ini sering keliru.
  \item \textbf{$\cos75^\circ=\cos45^\circ+\cos30^\circ$?} Fungsi trig tak
  menembus penjumlahan sudut.
\end{itemize}
\end{jebakan}

\begin{variasi}
$\cos15^\circ=\cos(45^\circ-30^\circ)=\frac{\sqrt6+\sqrt2}{4}$ (tanda plus
untuk selisih pada cos).
\end{variasi}

% ---------------- SOAL 4.3.2 ----------------
\soalno{4.3.2}{tan sudut campuran}
\begin{soalbox}{}
Hitung $\tan75^\circ$ dengan rumus jumlah sudut.
\end{soalbox}

\jawab
\langkah{1} $75^\circ=45^\circ+30^\circ$; $\tan45^\circ=1$,
$\tan30^\circ=\frac{1}{\sqrt3}$.
\langkah{2} $\tan(A+B)=\dfrac{\tan A+\tan B}{1-\tan A\tan B}$:
\[
\tan75^\circ=\frac{1+\frac{1}{\sqrt3}}{1-1\cdot\frac{1}{\sqrt3}}
=\frac{\sqrt3+1}{\sqrt3-1}\quad(\text{kali }\sqrt3\ \text{atas-bawah}).
\]
\langkah{3} Rasionalkan (kali $\sqrt3+1$):
\[
=\frac{(\sqrt3+1)^2}{(\sqrt3-1)(\sqrt3+1)}=\frac{3+2\sqrt3+1}{3-1}
=\frac{4+2\sqrt3}{2}=2+\sqrt3.
\]
\jadi $\tan75^\circ=2+\sqrt3\approx3{,}73$.

\begin{jebakan}
\begin{itemize}[leftmargin=*,topsep=2pt]
  \item \textbf{Lupa merasionalkan.} Bentuk $\frac{\sqrt3+1}{\sqrt3-1}$ belum
  ``rapi''; kalikan sekawan.
  \item \textbf{Tanda penyebut.} Untuk $\tan(A+B)$ penyebutnya $1-\tan A\tan B$
  (minus).
\end{itemize}
\end{jebakan}

\begin{variasi}
$\tan15^\circ=\tan(45^\circ-30^\circ)=2-\sqrt3$. Cek: $\tan75^\circ\cdot
\tan15^\circ=(2+\sqrt3)(2-\sqrt3)=1$ (berpenyiku!).
\end{variasi}

% ---------------- SOAL 4.3.3 ----------------
\soalno{4.3.3}{sudut rangkap dari cos}
\begin{soalbox}{}
Diketahui $\cos A=\tfrac{12}{13}$ ($A$ lancip). Tentukan $\sin2A$.
\end{soalbox}

\jawab
\langkah{1} Cari $\sin A$: dari segitiga $5$-$12$-$13$, $\sin A=\frac{5}{13}$
(positif, lancip).
\langkah{2} $\sin2A=2\sin A\cos A=2\cdot\dfrac{5}{13}\cdot\dfrac{12}{13}
=\dfrac{120}{169}$.
\jadi $\sin2A=\dfrac{120}{169}$.

\begin{jebakan}
\begin{itemize}[leftmargin=*,topsep=2pt]
  \item \textbf{$\sin2A=2\sin A$?} Salah; ada faktor $\cos A$ juga:
  $\sin2A=2\sin A\cos A$.
\end{itemize}
\end{jebakan}

\begin{variasi}
$\cos2A=2\cos^2 A-1=2\cdot\frac{144}{169}-1=\frac{119}{169}$; cek
$\sin^2 2A+\cos^2 2A=1$.
\end{variasi}

% ---------------- SOAL 4.3.4 ----------------
\soalno{4.3.4}{menurunkan sudut rangkap}
\begin{soalbox}{}
Buktikan $\sin2A=2\sin A\cos A$ dari rumus jumlah sudut.
\end{soalbox}

\jawab
\langkah{1} Tulis $2A=A+A$ dan pakai $\sin(A+B)=\sin A\cos B+\cos A\sin B$
dengan $B=A$:
\[
\sin(A+A)=\sin A\cos A+\cos A\sin A.
\]
\langkah{2} Kedua suku sama, jumlahkan: $=2\sin A\cos A$.
\jadi $\sin2A=2\sin A\cos A$. $\blacksquare$

\begin{jebakan}
\begin{itemize}[leftmargin=*,topsep=2pt]
  \item \textbf{Menganggap ini definisi.} Ia \emph{diturunkan} dari rumus jumlah,
  bukan aturan terpisah --- memahami asalnya membuatnya tak perlu dihafal.
\end{itemize}
\end{jebakan}

\begin{variasi}
Turunkan $\cos2A=\cos^2 A-\sin^2 A$ dari $\cos(A+A)$; lalu ubah ke
$1-2\sin^2 A$ memakai identitas pokok.
\end{variasi}

% ---------------- SOAL 4.3.5 ----------------
\soalno{4.3.5}{sin sudut setengah}
\begin{soalbox}{}
Hitung $\sin22{,}5^\circ$ dengan rumus sudut setengah.
\end{soalbox}

\jawab
\langkah{1} $22{,}5^\circ=\tfrac12(45^\circ)$, dan berada di kuadran I (positif).
\langkah{2} Rumus: $\sin\dfrac{\alpha}{2}=\sqrt{\dfrac{1-\cos\alpha}{2}}$ dengan
$\alpha=45^\circ$, $\cos45^\circ=\frac{\sqrt2}{2}$:
\[
\sin22{,}5^\circ=\sqrt{\frac{1-\frac{\sqrt2}{2}}{2}}
=\sqrt{\frac{2-\sqrt2}{4}}=\frac{1}{2}\sqrt{2-\sqrt2}.
\]
\jadi $\sin22{,}5^\circ=\dfrac12\sqrt{2-\sqrt2}\approx0{,}383$.

\begin{jebakan}
\begin{itemize}[leftmargin=*,topsep=2pt]
  \item \textbf{Salah tanda $\pm$.} Pilih $+$ karena $22{,}5^\circ$ di kuadran I.
  \item \textbf{Menyederhanakan pecahan bertingkat.} $\frac{1-\frac{\sqrt2}{2}}{2}
  =\frac{2-\sqrt2}{4}$ (kali $\frac22$ dulu di pembilang).
\end{itemize}
\end{jebakan}

\begin{variasi}
$\cos22{,}5^\circ=\frac12\sqrt{2+\sqrt2}$ (tanda $+$ di dalam akar). Cek
$\sin^2+\cos^2=1$.
\end{variasi}

% ---------------- SOAL 4.3.6 ----------------
\soalno{4.3.6}{mengenali sudut rangkap terbalik}
\begin{soalbox}{}
Sederhanakan $2\sin15^\circ\cos15^\circ$ menjadi nilai tunggal.
\end{soalbox}

\jawab
\langkah{1} Kenali pola $2\sin A\cos A=\sin2A$ dengan $A=15^\circ$:
\[
2\sin15^\circ\cos15^\circ=\sin(2\cdot15^\circ)=\sin30^\circ.
\]
\langkah{2} $\sin30^\circ=\dfrac12$.
\jadi $2\sin15^\circ\cos15^\circ=\dfrac12$.

\begin{jebakan}
\begin{itemize}[leftmargin=*,topsep=2pt]
  \item \textbf{Menghitung $\sin15^\circ$ dan $\cos15^\circ$ terpisah} lalu
  mengalikan --- boros dan rawan. Mengenali pola sudut rangkap jauh lebih cepat.
\end{itemize}
\end{jebakan}

\begin{variasi}
$\cos^2 A-\sin^2 A=\cos2A$; $1-2\sin^2 A=\cos2A$. Latih ``membaca mundur'' rumus
rangkap.
\end{variasi}

% ---------------- SOAL 4.3.7 ----------------
\soalno{4.3.7}{persamaan cos}
\begin{soalbox}{}
Selesaikan $\cos x=\dfrac{\sqrt3}{2}$ untuk $0^\circ\le x<360^\circ$.
\end{soalbox}

\jawab
\langkah{1} Sudut acuan: $\cos\alpha=\frac{\sqrt3}{2}\Rightarrow\alpha=30^\circ$.
\langkah{2} $\cos$ positif di kuadran I dan IV. Solusi:
\[
x=30^\circ \quad\text{atau}\quad x=360^\circ-30^\circ=330^\circ.
\]
\jadi $x\in\{30^\circ,\ 330^\circ\}$.

\begin{jebakan}
\begin{itemize}[leftmargin=*,topsep=2pt]
  \item \textbf{Hanya menulis $30^\circ$.} Dalam satu putaran, $\cos$ punya dua
  solusi (pasangan $\alpha$ dan $360^\circ-\alpha$).
\end{itemize}
\end{jebakan}

\begin{variasi}
$\cos x=-\frac12$ (KII \& KIII: $120^\circ,240^\circ$); $\cos x=0$
($90^\circ,270^\circ$).
\end{variasi}

% ---------------- SOAL 4.3.8 ----------------
\soalno{4.3.8}{persamaan tan}
\begin{soalbox}{}
Selesaikan $\tan x=1$ untuk $0^\circ\le x<360^\circ$.
\end{soalbox}

\jawab
\langkah{1} Sudut acuan: $\tan\alpha=1\Rightarrow\alpha=45^\circ$.
\langkah{2} $\tan$ berulang tiap $180^\circ$, jadi:
\[
x=45^\circ \quad\text{atau}\quad x=45^\circ+180^\circ=225^\circ.
\]
\jadi $x\in\{45^\circ,\ 225^\circ\}$.

\begin{jebakan}
\begin{itemize}[leftmargin=*,topsep=2pt]
  \item \textbf{Memakai $180^\circ-\alpha$ seperti sin.} Untuk $\tan$, solusi
  kedua adalah $\alpha+180^\circ$ (periode $180^\circ$, bukan $360^\circ$).
\end{itemize}
\end{jebakan}

\begin{variasi}
$\tan x=\sqrt3$ ($60^\circ,240^\circ$); $\tan x=-1$ ($135^\circ,315^\circ$).
\end{variasi}

% ---------------- SOAL 4.3.9 ----------------
\soalno{4.3.9}{persamaan sin linear}
\begin{soalbox}{}
Selesaikan $2\sin x-1=0$ untuk $0^\circ\le x<360^\circ$.
\end{soalbox}

\jawab
\langkah{1} Isolasi: $2\sin x=1\Rightarrow\sin x=\dfrac12$.
\langkah{2} Sudut acuan $\alpha=30^\circ$; $\sin$ positif di KI \& KII:
\[
x=30^\circ \quad\text{atau}\quad x=180^\circ-30^\circ=150^\circ.
\]
\jadi $x\in\{30^\circ,\ 150^\circ\}$.

\begin{jebakan}
\begin{itemize}[leftmargin=*,topsep=2pt]
  \item \textbf{Lupa mengisolasi $\sin x$ dulu.} Selesaikan aljabarnya
  ($\sin x=\frac12$) sebelum mencari sudut.
\end{itemize}
\end{jebakan}

\begin{variasi}
$2\cos x+\sqrt3=0\Rightarrow\cos x=-\frac{\sqrt3}{2}$; $\sqrt2\sin x=1$.
\end{variasi}

% ---------------- SOAL 4.3.10 ----------------
\soalno{4.3.10}{persamaan dengan sudut rangkap}
\begin{soalbox}{}
Selesaikan $\sin2x=\cos x$ untuk $0^\circ\le x<360^\circ$.
(Petunjuk: $\sin2x=2\sin x\cos x$.)
\end{soalbox}

\begin{ingat}{Jangan membagi dengan yang mungkin nol}
Godaan besar: membagi kedua ruas dengan $\cos x$. \emph{Jangan!} Itu membuang
solusi $\cos x=0$. Pindahkan semua ke satu ruas lalu \emph{faktorkan}.
\end{ingat}

\jawab
\langkah{1} Ganti $\sin2x=2\sin x\cos x$:
\[
2\sin x\cos x=\cos x.
\]
\langkah{2} Pindah ke satu ruas dan faktorkan (jangan dibagi!):
\[
2\sin x\cos x-\cos x=0 \Rightarrow \cos x\,(2\sin x-1)=0.
\]
\langkah{3} Dua kemungkinan:
\begin{itemize}[leftmargin=*,topsep=1pt]
  \item $\cos x=0 \Rightarrow x=90^\circ$ atau $270^\circ$;
  \item $2\sin x-1=0 \Rightarrow \sin x=\tfrac12 \Rightarrow x=30^\circ$ atau
  $150^\circ$.
\end{itemize}
\jadi $x\in\{30^\circ,\ 90^\circ,\ 150^\circ,\ 270^\circ\}$.

\begin{jebakan}
\begin{itemize}[leftmargin=*,topsep=2pt]
  \item \textbf{Membagi $\cos x$} sehingga solusi $90^\circ$ dan $270^\circ$
  hilang. Ini kesalahan HOTS paling umum --- \emph{faktorkan}, jangan bagi.
\end{itemize}
\end{jebakan}

\begin{variasi}
$\cos2x=\sin x$ (pakai $\cos2x=1-2\sin^2 x$, jadi kuadrat dalam $\sin x$);
$\sin2x=\sin x$.
\end{variasi}

% ============================================================
\begin{tips}
\textbf{Ringkasan strategi Bab 4.}
\begin{enumerate}[leftmargin=*,topsep=2pt]
  \item \emph{Pilih rasio} dari sisi diketahui \& dicari (SOH-CAH-TOA). Selalu
  gambar dan tandai sudut acuan.
  \item \emph{Sudut istimewa} berasal dari segitiga $45$-$45$-$90$ dan
  $30$-$60$-$90$ --- pahami, jangan hafal buta.
  \item \emph{Kuadran:} nilai mutlak dari sudut acuan, tanda dari ASTC.
  Identitas $\pm$ selalu dipilih tandanya lewat kuadran.
  \item \emph{Persamaan:} sin \& cos punya 2 solusi per putaran, tan berperiode
  $180^\circ$. \textbf{Faktorkan, jangan membagi} dengan fungsi yang bisa nol.
\end{enumerate}
\end{tips}

\end{document}
